\documentclass[conference]{IEEEtran}
\IEEEoverridecommandlockouts

% Essential IEEE Packages
\usepackage{cite}
\usepackage{amsmath,amssymb,amsfonts}
\usepackage{algorithmic}
\usepackage{graphicx}
\usepackage{textcomp}
\usepackage{xcolor}
\usepackage{booktabs}
\usepackage{url}
\usepackage{microtype}
\usepackage{tabularx}
\usepackage{array}
\usepackage{balance}

% Spacing parameters calibrated strictly for IEEE 6-page conference layout
\setlength{\textfloatsep}{7pt plus 1.5pt minus 1.5pt}
\setlength{\floatsep}{5pt plus 1.0pt minus 1.0pt}
\setlength{\intextsep}{5pt plus 1.0pt minus 1.0pt}
\setlength{\dbltextfloatsep}{7pt plus 1.5pt minus 1.5pt}

\begin{document}

\title{SAGRI: An Agro-Climatic Constraint-Gated Multimodal Architecture for Smallholder Advisory and Mandi Gate Settlement}

\author{
\IEEEauthorblockN{Shikhar Kesharwani\IEEEauthorrefmark{1}, Santusht Lakhanpal\IEEEauthorrefmark{1}, and Monu Singh\IEEEauthorrefmark{2}}
\IEEEauthorblockA{Department of Computer Science and Engineering, Bennett University, Greater Noida, UP, India\\
Email: \IEEEauthorrefmark{1}\{S24CSEU0502, S24CSEU0465\}@bennett.edu.in}
}

\maketitle

\begin{abstract}
Most digital agriculture tools built for Indian farmers fail when deployed in actual fields. Foliar disease models trained on laboratory images struggle with outdoor shadows and background soil. Crop recommendation models suggest water-thirsty crops in arid regions because they look only at soil nutrients. Meanwhile, market advisory tools stop at price predictions, growers often wait weeks for cash payouts at local mandis. We built SAGRI to address these issues together on low-cost hardware. To stop unrealistic crop advice, we placed ICAR agro-climatic rules as hard filters behind a Random Forest classifier. This step cut invalid crop suggestions from 15.8\% down to zero across test records. To estimate crop failure risk under drought, we joined XGBoost tree predictions with a five-variable physical crop-water deficit equation over 325,418 district records (1966--2017). This hybrid formula reached an ROC-AUC of 0.892. On the client side, our quantized MobileNetV2 ONNX model classifies 38 foliar disease states in 112~ms on a single CPU thread and takes 16.2~MB of storage. We also built an automated APMC settlement simulator that divides sale revenue into a 98\% farmer credit and a 2\% mandi cess. The system generates instant UPI transaction slips. All records persist locally in an embedded SQLite database running write-ahead logging. Queries execute in under 50~ms without relying on cloud infrastructure or steady mobile internet.
\end{abstract}

\begin{IEEEkeywords}
Agro-climatic constraints, foliar disease detection, MobileNetV2, XGBoost risk fusion, mandi price modeling, APMC settlement, offline-first computing.
\end{IEEEkeywords}

\section{Introduction}
India has roughly 146 million operational farm holdings, and 86\% of them measure under two hectares \cite{ref_nsso}. Smallholders on these plots rarely receive visits from trained agronomists. When crops develop leaf spots or fungal blight, farmers usually take a sample to a local agrochemical retailer. These shopkeepers often recommend expensive, broad-spectrum chemical sprays to maximize retail margins. The result is unnecessary debt for the farmer and long-term chemical damage to the soil.

Machine learning models could theoretically offer low-cost advice on basic Android phones. In practice, models from research papers often make dangerous mistakes in the field. First, standard crop recommendation classifiers optimize purely mathematical loss functions. If a Random Forest model sees high soil potassium and nitrogen from a district in western Rajasthan, it frequently recommends wetland paddy (\textit{Oryza sativa}). The model does not know that the district receives only 350~mm of annual rainfall and that local groundwater is already depleted. Second, many commercial agritech apps artificially push low model confidence scores into the 90\% to 99\% range. This artificial inflation reassures the user, but it hides genuine diagnostic uncertainty. Third, advisory tools rarely connect to market cash flows. Farmers still face delayed cash payments from commission agents (\textit{arhtiyas}) at Agricultural Produce Market Committee (APMC) mandis, Growers often wait two to six weeks before receiving cash.

We designed SAGRI (\textit{Smart Agriculture Real-time Intelligence}) to solve these problems in an integrated platform. We placed an ICAR-based biological constraint layer behind our crop recommendation model. The layer prunes out crops that cannot survive local thermal and rainfall conditions before any ranking is shown to the user. For climate failure risk, we built a hybrid radar that combines statistical predictions from XGBoost with a deterministic crop-water deficit formula. To assist with mandi transactions, we added an APMC settlement simulator that calculates statutory market fees (98\% farmer credit and 2\% mandi cess) and creates immediate UPI transaction slips. Finally, we stored all data in a local SQLite database running write-ahead logging (WAL). Field agents can query logs and record predictions offline on inexpensive laptops without paying for cloud servers or waiting for intermittent cellular connections.

\section{Related Work}

\subsection{Mobile Plant Pathology on Budget Hardware}
Deep learning became the primary approach for plant pathology after Mohanty \textit{et al.} \cite{ref_mohanty} reached high diagnostic scores on the PlantVillage dataset. However, models that succeed in laboratories often struggle in open fields. Ferentinos \cite{ref_ferentinos} evaluated deep architectures like VGG-16 \cite{ref_vgg} and GoogLeNet, He showed that large convolutional filters detect foliar spots well, but they take up over 500~MB of memory. ResNet-50 \cite{ref_resnet} improved gradient backpropagation through residual shortcuts, but its 98~MB file size still causes noticeable thermal throttling on low-end Android phones. Sandler \textit{et al.} introduced MobileNetV2 \cite{ref_mobilenet} to reduce parameters through inverted residual bottlenecks. In this project, we quantize MobileNetV2 to 8-bit integers and run it via ONNX Runtime. The resulting model requires only 16.2~MB of memory and finishes inference in 112~ms on a single CPU thread without GPU support.

\subsection{Agronomic Constraints in Crop Selection}
Tabular crop recommendation has traditionally relied on tree ensembles and support vector machines. Medar \textit{et al.} \cite{ref_medar} showed that Random Forest \cite{ref_rf} outperforms simple decision trees and SVMs \cite{ref_svm} when predicting crop suitability from soil nutrients. However, standard classifiers treat soil chemistry and climate features as unconstrained numerical inputs. If a training set contains nutrient combinations common to both rice and pulses, the classifier will predict rice in dry zones whenever the NPK values match. While constraint satisfaction methods \cite{ref_tsang} are common in operations research, agricultural machine learning papers rarely enforce biological boundaries over statistical outputs. SAGRI applies these boundaries directly through ICAR regional exclusion rules.

\subsection{Climate Stress Metrics and Mandi Price Forecasting}
Crop yield vulnerability under extreme weather has received substantial attention. Lobell \textit{et al.} \cite{ref_lobell} showed how heat spikes reduce cereal harvests across multiple seasons. Challinor \textit{et al.} \cite{ref_challinor} demonstrated that physical evapotranspiration shortfalls capture crop failure dynamics better than raw temperature averages alone. Gradient-boosted decision trees such as XGBoost \cite{ref_xgboost} handle tabular climate data effectively, but they can produce unexpected predictions during severe weather outliers that did not appear in training data. Combining tree predictions with physical crop-water deficit math prevents models from underestimating risk during intense droughts.
On the commodity pricing side, earlier research by Kumar \textit{et al.} \cite{ref_price_arima} applied ARIMA models to wholesale Agmarknet records. Jha and Sinha \cite{ref_jha} observed that raw price models often confuse broad monetary inflation with genuine crop price increases. Building on their findings, we deflate 25 years of Agmarknet historical figures with official Consumer Food Price Index (CFPI) data and feed real-time weather anomalies from Open-Meteo \cite{ref_openmeteo} into our price forecasting pipeline.

% Full-width architectural diagram across both columns
\begin{figure*}[t]
\centering
\resizebox{0.96\textwidth}{!}{
\begin{tabular}{|c|c|c|}
\hline
\rule{0pt}{14pt}\textbf{Client Experience Layer} & \textbf{Unified Inference Core (FastAPI :8000)} & \textbf{Persistence \& Settlement Rails} \\[3pt]
\hline
\rule{0pt}{26pt}\begin{tabular}{c}
\textbf{SAGRI Classic (Port 5173)} \\
React 18 $\cdot$ Tailwind $\cdot$ Voice Assistant \\[3pt]
\hline
\rule{0pt}{14pt}\textbf{SAGRI Modern (Port 5174)} \\
Radix UI $\cdot$ Command Palette (\texttt{Ctrl+K})
\end{tabular}
&
\begin{tabular}{c}
\textbf{MobileNetV2 ONNX} (16.2 MB, 112 ms CPU) \\
\textbf{Constraint-Gated RF} (ICAR Bioclimatic Filter) \\
\textbf{Hybrid Risk Radar} (XGBoost + Shortfall Engine) \\
\textbf{Price Forecaster} (25-yr CFPI + Open-Meteo)
\end{tabular}
&
\begin{tabular}{c}
\textbf{SQLite WAL Engine (\texttt{sagri.db})} \\
Zero-cloud offline-first ACID storage \\[3pt]
\hline
\rule{0pt}{14pt}\textbf{APMC 98/2 Split Settlement} \\
UPI rail simulation with instant UTR
\end{tabular}
\\[18pt]
\hline
\end{tabular}
}
\caption{System topology: concurrent client interfaces communicating with the multi-model FastAPI service and local SQLite storage layer.}
\label{fig:arch}
\end{figure*}

\section{System Architecture}
Fig.~\ref{fig:arch} shows the overall software architecture. Rather than relying on a single interface, we built two separate frontend clients to match different field needs.
The first client is SAGRI Classic (Port 5173). It uses React 18 with large buttons, high-contrast colors, spoken prompts, and simple WebGL soil profile graphics. This view serves field extension workers and farmers who prefer voice guidance and clear, simple layouts.
The second client is SAGRI Modern (Port 5174). It is a dense workstation interface built with Radix UI, keyboard shortcuts (\texttt{Ctrl+K} command palette), and detailed statistical graphs. This view serves agronomists, researchers, and APMC market secretaries who analyze records in bulk.

Both client applications communicate over REST endpoints with a local FastAPI backend running on Port 8000. For data persistence, we avoided remote cloud databases that become inaccessible during rural network outages. Instead, we use an embedded SQLite database file (\texttt{sagri.db}). Setting \texttt{PRAGMA journal\_mode = WAL} enables concurrent reads and writes, so both frontend applications can log predictions and inspect history simultaneously without database lockups.

\section{Data Curation and Preprocessing}
Table~\ref{tab:datasets} summarizes the four datasets and partition schemes used across our training pipelines.

\begin{table}[htbp]
\caption{Dataset Summary and Partitioning Strategy}
\label{tab:datasets}
\centering
\resizebox{\columnwidth}{!}{
\begin{tabular}{llrl}
\toprule
\textbf{Dataset Source} & \textbf{Target Task} & \textbf{Sample Count} & \textbf{Evaluation Split} \\
\midrule
PlantVillage & Foliar Pathology & 54,305 Images & 80\% Train / 20\% Test \\
ICAR Soil Survey & Crop Recommendation & 2,200 Rows & 5-Fold Stratified CV \\
ICRISAT + IMD + ERA5 & Failure Risk Radar & 325,418 Rows & 1966--2007 Train / 2008--2017 Test \\
Agmarknet 25-Year & Mandi Forecasting & 1.2M Records & Chronological 80/20 Holdout \\
\bottomrule
\end{tabular}
}
\end{table}

Preparing real-world data required substantial preprocessing. For the 54,305 PlantVillage leaf images, we removed corrupted files and out-of-focus samples, then cropped foliage into standard 224$\times$224 boxes. For the 2,200 crop recommendation records, we inspected distributions and clamped extreme soil pH values that exceeded three standard deviations.

Constructing the climate failure radar required extensive data engineering. We merged 52 years (1966--2017) of ICRISAT district yield data with daily rainfall grids from the India Meteorological Department (IMD), ERA5 surface wind speeds, and solar radiation records from NASA POWER. Because dozens of Indian districts split or shifted administrative borders during those five decades, we mapped historical boundaries onto Census 2011 district identifiers to keep spatial comparisons consistent. We defined a crop failure ($Y=1$) as any season where yield dropped more than 20\% below the district's rolling 10-year mean. For mandi price forecasting, we collected 25 years of daily modal price records across 30 crops from Agmarknet. We matched each record to monthly national food inflation tables.

\section{Methodology and Model Design}

\subsection{Pipeline 1: MobileNetV2 Foliar Pathology Scanner}
We trained our foliar diagnostic network on 38 crop-disease classes (covering 14 hosts) using the PlantVillage repository. The model uses inverted residual blocks with linear bottlenecks. For each incoming image $\mathbf{I}$, we apply center cropping and standard ImageNet normalization:
\begin{equation}
\mathbf{I}_{\text{crop}} = \text{CenterCrop}(\mathbf{I}, \min(H, W))
\end{equation}
\begin{equation}
\mathbf{I}_{\text{norm}} = \frac{\text{Resize}(\mathbf{I}_{\text{crop}}, 224, 224) - \boldsymbol{\mu}}{\boldsymbol{\sigma}}
\end{equation}
where channel means are $\boldsymbol{\mu} = [0.485, 0.456, 0.406]$ and standard deviations are $\boldsymbol{\sigma} = [0.229, 0.224, 0.225]$. Outdoor lighting varies substantially across farm fields. To stabilize predictions against shadows and slight camera tilts, we run two-pass Test-Time Augmentation (TTA). The model averages outputs from the upright image and a horizontal flip:
\begin{equation}
\mathbf{z} = \frac{\text{softmax}(f_\theta(\mathbf{I}_{\text{norm}})) + \text{softmax}(f_\theta(\text{Flip}_H(\mathbf{I}_{\text{norm}})))}{2}
\end{equation}
The output vector $\mathbf{z}$ displays raw probabilities directly to the user. We do not use artificial confidence thresholds or score boosts. If confidence is low, the platform displays the three closest candidate diagnoses alongside chemical options registered with the Central Insecticides Board \& Registration Committee (CIBRC) and verified organic treatments.

\subsection{Pipeline 2: Agronomic Constraint-Gated Crop Recommender}
We trained an ensemble of 100 decision trees on 2,200 ICAR soil records across 22 Indian crops. Each input row contains seven parameters:
\begin{equation}
\mathbf{x} = [N, P, K, \text{pH}, T_{\text{mean}}, H_{\text{rel}}, R_{\text{ann}}]
\end{equation}
Instead of directly returning the highest raw probability output $\mathbf{p} = [p_1, p_2, \ldots, p_{22}]$, our inference pipeline passes candidate recommendations through two biological filters.

\textit{Stage 1 (Regional Incompatibility Filter):} If crop $c$ is listed under ICAR exclusion records $\mathcal{E}(r)$ for agro-climatic zone $r$, its probability becomes zero:
\begin{equation}
p'_c = 0 \quad \forall \ c \in \mathcal{E}(r)
\end{equation}

\textit{Stage 2 (Physiological Threshold Gate):} If seasonal rainfall is below the biological survival minimum $W_{\min}(c)$ or temperatures fall outside safe reproductive ranges $[T_{\min}(c), T_{\max}(c)]$, the crop is removed:
\begin{equation}
p''_c = \begin{cases} 
0, & \text{if } R < W_{\min}(c) \text{ or } T \notin [T_{\min}, T_{\max}] \\ 
p'_c, & \text{otherwise} 
\end{cases}
\end{equation}
The surviving probabilities are re-normalized:
\begin{equation}
\mathbf{p}^* = \frac{\mathbf{p}''}{\sum_k p''_k}
\end{equation}
This two-stage check ensures that mathematical correlation does not override basic plant biology.

\subsection{Pipeline 3: Hybrid Risk Radar}
To predict crop failure, we joined 325,418 district-year records from ICRISAT (1966--2017) with daily IMD rainfall grids, ERA5 wind measurements, and NASA POWER solar data. A failure label ($Y=1$) marks seasons where yield fell more than 20\% below the district's 10-year rolling mean.

Our XGBoost model estimates statistical failure probability $P_{\text{XGB}} \in [0, 1]$ over 15 tabular features. In parallel, our physics engine computes an environmental stress shortfall $S \in [0, 100]$ across five direct crop stresses:
\begin{align}
S_{\text{heat}} &= \max\left(0, \frac{T_{\max} - T_{\text{crit}}(c)}{10}\right) \times 20 \\
S_{\text{rain}} &= \max\left(0, \frac{W_{\text{req}}(c) - R_{\text{obs}}}{W_{\text{req}}(c)}\right) \times 25 \\
S_{\text{et}}   &= \max\left(0, \frac{E_T - R_{\text{obs}}}{E_T}\right) \times 20 \\
S_{\text{npk}}  &= \max\left(0, \frac{\|\mathbf{N}_{\text{req}} - \mathbf{N}_{\text{obs}}\|_1}{\|\mathbf{N}_{\text{req}}\|_1}\right) \times 20 \\
S_{\text{wind}} &= \max\left(0, v_{\text{wind}} - 8.0\right) \times 15
\end{align}
\begin{equation}
S_{\text{total}} = \min(100, S_{\text{heat}} + S_{\text{rain}} + S_{\text{et}} + S_{\text{npk}} + S_{\text{wind}})
\end{equation}
The combined risk score balances both components:
\begin{equation}
\text{Risk} = \left(0.65 \times P_{\text{XGB}} + 0.35 \times \frac{S_{\text{total}}}{100}\right) \times 100\%
\end{equation}
We set the 0.65/0.35 weights through grid search on held-out validation districts. This combination prevents the model from underestimating risk during intense droughts that had no historical counterpart in the training records.

\subsection{Pipeline 4: Inflation-Adjusted Price Forecaster}
Nominal crop prices in India are distorted by decades of rupee inflation. We normalized 25 years of daily modal price records across 30 crops using monthly Consumer Food Price Index (CFPI) numbers:
\begin{equation}
P_{\text{real}, t} = P_{\text{nominal}, t} \times \left(\frac{\text{CFPI}_{\text{base}}}{\text{CFPI}_t}\right)
\end{equation}
We then fit an autoregressive time-series model of order $p$ on the deflated series, The regression incorporates live temperature and precipitation anomalies $\mathbf{W}_t$ pulled dynamically from Open-Meteo:
\begin{equation}
P_{\text{real}, t} = \phi_0 + \sum_{i=1}^p \phi_i P_{\text{real}, t-i} + \boldsymbol{\theta}^T \mathbf{W}_t + \epsilon_t
\end{equation}
This keeps price forecasts tied to actual purchasing power changes rather than general monetary inflation.

\subsection{Pipeline 5: Multi-Persona Conversational Advisory}
We configured Google Gemini 1.5 Flash with four specialized agronomic prompts: General Extension Advisor, Soil Chemistry Specialist, Mandi Market Analyst, and Integrated Pest Management (IPM) Pathologist. The system supports natural dialogue in Hindi, Punjabi, and English. An automated Unicode check detects the user's script and sets the response language accordingly.

\section{APMC Mandi Settlement and Data Ethics}
Delayed payment remains a persistent hardship for Indian smallholders. A farmer brings produce to a regulated APMC yard, agrees to an auction price, and unloads the harvest. However, cash payments from commission agents (\textit{arhtiyas}) often stall for weeks while paperwork circulates among traders. Farmers frequently borrow high-interest cash just to buy seed and fertilizer for the following season.

To show how digital advisory tools can address this problem, SAGRI integrates payment settlement directly into the market advisory interface. When a farmer enters an expected sale of $Q$ quintals at forecasted price $P$, total proceeds are computed:
\begin{equation}
\text{Gross} = P \times Q
\end{equation}
The platform divides these proceeds under statutory APMC rules:
\begin{equation}
\text{Farmer Net} = 0.98 \times \text{Gross}, \quad \text{Mandi Cess} = 0.02 \times \text{Gross}
\end{equation}
The 98\% share simulates an immediate direct payout to the farmer's Virtual Payment Address (VPA) over UPI rails, while the 2\% statutory fee routes to the market committee account. The system generates a digital receipt with a 12-digit Unique Transaction Reference (UTR), timestamp, and fee summary.

Data privacy is equally essential. To comply with India's Digital Personal Data Protection (DPDP) Act \cite{ref_dpdp}, all farmer records, diagnostic photos, and financial entries stay inside the local SQLite database. No parcel boundary coordinates, contact numbers, or transaction logs are exported to commercial advertisers or commodity trading firms.

\section{Experimental Results}

\subsection{Execution Latency and Pipeline Accuracy}
We benchmarked the complete software stack on an unaccelerated quad-core Intel Core i5-1135G7 CPU at 2.40~GHz with 8~GB of RAM (Table~\ref{tab:benchmarks}). This setup reflects budget laptops used by rural extension workers and primary agricultural societies. The quantized MobileNetV2 ONNX model averaged 112~ms per image on a single CPU thread and consumed 16.2~MB of RAM. Two-pass horizontal flip TTA improved field diagnostic accuracy by 1.6\% without GPU requirements. SQLite with WAL mode sustained 1,840 write operations per second. This setup prevented database stalls during busy morning mandi weigh-in hours.

\begin{table}[htbp]
\caption{System Performance Across Pipeline Modules}
\label{tab:benchmarks}
\centering
\resizebox{\columnwidth}{!}{
\begin{tabular}{llr}
\toprule
\textbf{Pipeline Component} & \textbf{Evaluation Metric} & \textbf{Observed Value} \\
\midrule
Foliar Scanner (ONNX) & Top-1 Field Accuracy & 73.4\% \\
Foliar Scanner (ONNX) & Single-Thread CPU Latency & 112 ms \\
Foliar Scanner (ONNX) & Model File Size & 16.2 MB \\
Crop Recommender & 22-Crop Test Accuracy & 99.3\% \\
Risk Radar (Hybrid) & Out-of-Time ROC-AUC & 0.892 \\
Mandi Price Forecaster & 30-Day Mean Abs. Pct. Error & 8.35\% \\
Backend Service (FastAPI) & Median (P50) API Response & 48 ms \\
Storage Engine (SQLite WAL) & Write Throughput & 1,840 tx/s \\
\bottomrule
\end{tabular}
}
\end{table}

\subsection{Comparative Baseline Evaluation}
To see where lighter models succeed and where heavier architectures waste compute resources, we benchmarked candidate models across identical data splits (Table~\ref{tab:baselines}).

\begin{table}[htbp]
\caption{Model Comparisons Across Identical Partitions}
\label{tab:baselines}
\centering
\resizebox{\columnwidth}{!}{
\begin{tabular}{llrrr}
\toprule
\textbf{Task} & \textbf{Architecture} & \textbf{Acc. / AUC} & \textbf{CPU Latency} & \textbf{RAM} \\
\midrule
Leaf Diagnostics & VGG-16 \cite{ref_vgg} & 72.8\% & 980 ms & 528 MB \\
Leaf Diagnostics & ResNet-50 \cite{ref_resnet} & 74.1\% & 410 ms & 98 MB \\
Leaf Diagnostics & \textbf{MobileNetV2 (Ours)} & \textbf{73.4\%} & \textbf{112 ms} & \textbf{16.2 MB} \\
\midrule
Crop Choice & Decision Tree (CART) & 84.1\% & 1 ms & 0.4 MB \\
Crop Choice & Support Vector Machine & 89.2\% & 14 ms & 3.8 MB \\
Crop Choice & \textbf{Random Forest (Ours)} & \textbf{99.3\%} & \textbf{6 ms} & \textbf{2.1 MB} \\
\midrule
Yield Failure Risk & Logistic Regression & 0.712 & 2 ms & 0.1 MB \\
Yield Failure Risk & Random Forest Regressor & 0.835 & 12 ms & 8.4 MB \\
Yield Failure Risk & Standalone XGBoost & 0.871 & 5 ms & 4.2 MB \\
Yield Failure Risk & \textbf{Hybrid Radar (Ours)} & \textbf{0.892} & \textbf{8 ms} & \textbf{4.3 MB} \\
\bottomrule
\end{tabular}
}
\end{table}

In foliar diagnostics, ResNet-50 achieved 74.1\% field accuracy. That is only 0.7 percentage points higher than our quantized MobileNetV2 (73.4\%). However, ResNet-50 required 410~ms per image and 98~MB of storage. On entry-level Android smartphones (such as MediaTek Helio G25 or Snapdragon 439 devices) used by field staff, this delay freezes the camera viewfinder. VGG-16 performed much worse: it used 528~MB of memory and took nearly a full second (980~ms) per scan. Quantized MobileNetV2 ran in 112~ms with a 16.2~MB footprint. This speed gave field personnel immediate diagnostic feedback.

In crop selection, simple CART decision trees (84.1\%) and SVMs (89.2\%) struggled with non-linear relationships between soil NPK chemistry and rainfall. They trailed our Random Forest model (99.3\%) by more than ten percentage points.

For climate risk, logistic regression performed poorly with an ROC-AUC of 0.712. Standalone XGBoost achieved 0.871 AUC, but it missed yield crashes caused by sudden localized heatwaves. Fusing XGBoost with our physical crop-water shortfall formula raised ROC-AUC to 0.892 (+0.021 gain). The physical formula caught drought stress patterns that tree splits missed.

\subsection{Ablation Study Across System Modules}
We removed components one at a time to examine whether each part provides measurable value (Table~\ref{tab:ablation}).

\begin{table}[htbp]
\caption{Systematic Ablation Across Architecture Modules}
\label{tab:ablation}
\centering
\resizebox{\columnwidth}{!}{
\begin{tabular}{llr}
\toprule
\textbf{Pipeline Component} & \textbf{Ablation Variant} & \textbf{Observed Impact / Failure Mode} \\
\midrule
Full Architecture & Default SAGRI Stack & 0.892 AUC, 8.35\% MAPE, 0\% violations \\
Crop Engine & w/o ICAR Constraint Mask & 15.8\% unviable crops advised in arid zones \\
Climate Risk & w/o Crop-Water Shortfall & ROC-AUC drops from 0.892 to 0.871 in heatwaves \\
Mandi Pricing & w/o CFPI Inflation Deflation & 30-day forecast error jumps to 14.2\% MAPE \\
Leaf Vision & w/o Two-Pass Augmentation & Field top-1 accuracy falls from 73.4\% to 71.8\% \\
Local Storage & w/o SQLite WAL Mode & Concurrency lockups appear at $\ge$15 sessions \\
\bottomrule
\end{tabular}
}
\end{table}

Removing the ICAR biological constraint mask led to immediate agronomic errors. The unconstrained Random Forest advised water-intensive crops such as paddy and sugarcane in arid western Rajasthan districts in approximately 16 out of every 100 test queries. In a real farming scenario, farmers following these unconstrained predictions would deplete borewells and suffer severe financial losses. Removing the physical crop-water shortfall formula caused risk ROC-AUC to drop from 0.892 to 0.871. This confirms that tree models struggle when temperatures exceed historical records. Without CFPI price deflation, forecast error jumped from 8.35\% to 14.2\% MAPE. The model confused general currency inflation with real commodity price shifts. Finally, SQLite without WAL mode caused database lockups once concurrent API requests reached 15 sessions.

\subsection{Known Limitations and Field Edge Cases}
Stress testing uncovered an important limitation tied to the PlantVillage dataset: the dataset has no diseased blueberry leaves. It contains only the class \texttt{Blueberry\_\_\_healthy}. When we tested our ONNX model on field images of nutrient-starved and sunburned blueberry leaves from local nurseries, the model had no matching disease category. It classified the leaves as \texttt{Blueberry\_\_\_healthy}, but with a low softmax confidence score of 26.3\%.

Commercial agriculture tools often mask low confidence by artificially scaling scores up to 90\%. In contrast, our platform presents raw probabilities directly. This transparency makes the model's uncertainty visible to the user. To prevent misleading negative results, we added a rule: whenever the model predicts a healthy leaf with under 45\% confidence, the system flags the result as ``Inconclusive / Suspected Abiotic Stress''. This alert encourages farmers to request a manual check from a Krishi Vigyan Kendra (KVK) officer rather than assuming the plant is unharmed.

\section{Conclusion and Future Scope}
Machine learning models often make impractical recommendations in agriculture when they operate without physical domain constraints. They may suggest water-hungry crops in arid zones, hide diagnostic uncertainty behind inflated confidence scores, or mistake general inflation for market price trends. In this paper, we presented SAGRI, an architecture that grounds agricultural algorithms in physical constraints and market systems. We placed ICAR biological rules behind crop recommendation trees, paired tree predictions with physical crop-water deficit math, deflated historical mandi prices against 25 years of food inflation, and simulated automated UPI payments. The resulting platform gives smallholders practical, calibrated advice. Our local SQLite database eliminates cloud dependencies and keeps farmer data on-premise under DPDP regulations. In future work, we plan to expand the biological constraint tables across all 128 Indian agro-ecological zones and evaluate portable micro-spectrometers for direct soil nutrient measurement.

\balance

\begin{thebibliography}{00}
\bibitem{ref_nsso} National Sample Survey Office, ``Situation Assessment of Agricultural Households in India,'' Ministry of Statistics and Programme Implementation, New Delhi, Tech. Rep. 587, 2019.
\bibitem{ref_mohanty} S. P. Mohanty, D. P. Hughes, and M. Salath\'e, ``Using deep learning for image-based plant disease detection,'' \textit{Frontiers in Plant Science}, vol. 7, p. 1419, Sep. 2016.
\bibitem{ref_ferentinos} K. P. Ferentinos, ``Deep learning models for plant disease detection and diagnosis,'' \textit{Computers and Electronics in Agriculture}, vol. 145, pp. 311--318, Feb. 2018.
\bibitem{ref_vgg} K. Simonyan and A. Zisserman, ``Very deep convolutional networks for large-scale image recognition,'' in \textit{Proc. Int. Conf. Learn. Represent. (ICLR)}, 2015.
\bibitem{ref_resnet} K. He, X. Zhang, S. Ren, and J. Sun, ``Deep residual learning for image recognition,'' in \textit{Proc. IEEE Conf. Comput. Vis. Pattern Recognit. (CVPR)}, 2016, pp. 770--778.
\bibitem{ref_mobilenet} M. Sandler, A. Howard, M. Zhu, A. Zhmoginov, and L.-C. Chen, ``MobileNetV2: Inverted residuals and linear bottlenecks,'' in \textit{Proc. IEEE/CVF Conf. Comput. Vis. Pattern Recognit. (CVPR)}, 2018, pp. 4510--4520.
\bibitem{ref_rf} L. Breiman, ``Random forests,'' \textit{Machine Learning}, vol. 45, no. 1, pp. 5--32, 2001.
\bibitem{ref_svm} C. Cortes and V. Vapnik, ``Support-vector networks,'' \textit{Machine Learning}, vol. 20, no. 3, pp. 273--297, 1995.
\bibitem{ref_medar} R. J. Medar, V. S. Rajpurohit, and S. Shweta, ``Crop yield prediction using machine learning techniques,'' in \textit{Proc. IEEE Int. Conf. Comput. Commun. (ICCMC)}, 2019, pp. 906--910.
\bibitem{ref_tsang} E. Tsang, \textit{Foundations of Constraint Satisfaction}. London, UK: Academic Press, 1993.
\bibitem{ref_lobell} D. B. Lobell, W. Schlenker, and J. Costa-Roberts, ``Climate trends and global crop production since 1980,'' \textit{Science}, vol. 333, no. 6042, pp. 616--620, Jul. 2011.
\bibitem{ref_challinor} A. J. Challinor, T. R. Wheeler, C. Garforth, P. Craufurd, and A. Kassam, ``Assessing the vulnerability of food crop systems in Africa to climate change,'' \textit{Climatic Change}, vol. 83, no. 3, pp. 381--399, Nov. 2007.
\bibitem{ref_xgboost} T. Chen and C. Guestrin, ``XGBoost: A scalable tree boosting system,'' in \textit{Proc. 22nd ACM SIGKDD Int. Conf. Knowl. Discov. Data Min.}, 2016, pp. 785--794.
\bibitem{ref_price_arima} S. Kumar, P. Sharma, and A. Patel, ``Price forecasting of commodities using machine learning techniques,'' \textit{Int. J. Comput. Appl.}, vol. 183, no. 15, pp. 35--40, 2021.
\bibitem{ref_jha} G. R. Jha and S. Sinha, ``Agricultural price forecasting using neural network model: An innovative information delivery system,'' \textit{Agric. Econ. Res. Rev.}, vol. 26, no. 2, pp. 229--239, 2013.
\bibitem{ref_openmeteo} P. Zippenfenig, ``Open-Meteo: Free weather API for open-source research,'' 2023. [Online]. Available: https://open-meteo.com/
\bibitem{ref_nabard} National Bank for Agriculture and Rural Development, ``Kisan Credit Card Scheme Guidelines,'' NABARD, Mumbai, Cir. 12/2023, 2023.
\bibitem{ref_dpdp} Ministry of Law and Justice, ``The Digital Personal Data Protection Act, 2023,'' \textit{The Gazette of India}, Act No. 22 of 2023, Aug. 2023.
\end{thebibliography}

\end{document}
